%% sec7_4.tex — Section 7.4: Exercises
%% Chapter 7 of "A First Course in Monte Carlo Methods"
%% Sanz-Alonso & Al-Ghattas, University of Chicago

\section{Exercises}
\label{sec:7.4}

\begin{enumerate}

\item[7.1] Consider the function
\[
  V(x) = \Bigl(\cos(50x) + \sin(20x)\Bigr)^2
\]
defined for $0 < x < 1$.
\begin{enumerate}
  \item Use the command \texttt{scipy.optimize} to identify a maximum at $x^* = 0.379$ with a
    value of $V(x^*) = 3.8325$.
  \item Design a simulated annealing algorithm to optimize the function $h$. You may build on
    the following sketch or write your code from scratch:
\begin{verbatim}
diff = iter = 1
while diff > 10**(-4):
    prop = x[iter-1] + random.uniform(-1, 1)*scale
    u = random.uniform(0, 1)
    if not (0 <= prop <= 1) or
            (math.log(u)*temp[iter-1] > V(prop) - V(x[iter-1])):
        prop = x[iter-1]
    x.append(prop)
    Vcur = V(prop)
    Vval.append(Vcur)
    if iter > 10 and len(set(x[(iter//2):iter])) > 1:
        diff = max(Vval) - max(Vval[0:(iter//2)])
    iter += 1
\end{verbatim}
    This algorithm suggests stopping when the observed maximum has not changed in the second
    half of the sequence $\{x_t\}$. The constraint involving \texttt{unique} is canceling the
    stopping rule when no perturbation has been accepted in the second half of the iterations,
    meaning that the scale may be inappropriate. (Note: the updates of \texttt{temp} and
    \texttt{scale} need to be included in the loop.)

    Consider the following choices:\\
    (i) A scale defined by $\sqrt{T_t}$ and a temperature $T_t$ decrease in $1/\log(1+t)$.\\
    (ii) A scale defined by $5\sqrt{T_t}$ and a temperature decrease in $1/(1+t)^2$.\\
    (iii) Decrease the scale in (i) and (ii) by a factor of 10. What do you observe?
\end{enumerate}

\item[7.2] In this exercise, you will implement and compare the MALA and parallel tempering for
sampling from a double well potential $V(x) = 2(x^2 - 1)^2$.
\begin{enumerate}
  \item Implement MALA to sample from $V(x)$, generating $10^5$ samples with a fixed
    step-size $\epsilon = 0.05$. Plot the trajectory of the samples.
  \item Implement parallel tempering using 4 chains at temperatures $1, 1/2, 1/4, 1/8$ with a
    fixed step-size $\epsilon = 0.05$, generating $10^5$ samples per chain. After every 10
    samples, randomly select two chains and attempt a swap using the Metropolis criterion. Plot
    the trajectories for all chains.
\end{enumerate}
\begin{enumerate}
  \setcounter{enumii}{0}
  \item[2.] \begin{enumerate}
    \item For MALA, perform 10 runs to sample from $V(x)$, recording the number of steps
      required for transitions from $x < 0$ to $x > 0$ and vice versa up to $10^5$ samples.
      Assume the starting position is randomly chosen from the standard normal distribution.
      Compute the average transition steps over the runs.
    \item Repeat the above for parallel tempering.
    \item Compare the average transition times for MALA and parallel tempering, discussing the
      efficiency and computational cost of each method.
  \end{enumerate}
  \item[3.] An alternative method for state swapping in parallel tempering is to attempt to swap
    between adjacent chains sequentially. Evaluate the impact of this method on the algorithm's
    efficiency.
\end{enumerate}

\item[7.3] Consider the Rosenbrock function $f : \R^2 \to \R$ defined by
\[
  f(x, y) = 100(y - x^2)^2 + (1 - x)^2.
\]
In this question, you will implement and compare simulated annealing, gradient descent and
Newton descent to minimize this function. Writing $z = (x, y)$, recall that in gradient
descent, we iterate:
\[
  z^{(k)} = z^{(k-1)} - \eta \nabla f(z^{(k-1)}),
\]
and in Newton descent, we iterate:
\[
  z^{(k)} = z^{(k-1)} - \eta\bigl(\nabla^2 f(z^{(k-1)})\bigr)^{-1} \nabla f(z^{(k-1)}),
\]
where $\eta > 0$ is a step-size parameter. Using a starting point of $z^{(0)} = (-1.2, 1)$ and
a step-size of your choosing, run each of the algorithms until you reach the termination
condition $\|\nabla f(z^{(k)})\|_2 \leq 10^{-6}$. Report the obtained solution from each
scheme, as well as the number of iterations required. Provide a plot of the iterates from the
three schemes super-imposed on a plot of the function. You may use a 3d plot or a 2d contour
plot.

\item[7.4] Consider the function $V : \R \to \R$, defined by
\[
  V(x; c) = \frac{1}{2}(x^2 - c)^2,
\]
and where $c > 0$ is a parameter controlling the bi-modality of $V$. In this question, you will
compare the efficiency of MALA and parallel tempering applied to $V$ for various choices of
parameter $c \in \{0, 1, 5\}$.
\begin{enumerate}
  \item Implement MALA with $N = 10^4$ and $\epsilon = 0.1$.
  \item Implement parallel tempering with $N = 10^4$, attempting swaps after every 10 samples
    and using the following temperatures $T \in \{0.01, 0.1, 0.5, 1\}$.
  \item Compare the two algorithms by plotting a histogram and path of the samples for each
    value of $c$. Explain how (and why) increasing $c$ affects the relative performance of the
    two algorithms.
\end{enumerate}

\end{enumerate}
